% SPDX-License-Identifier: CC-BY-SA-4.0
% Author: Matthieu Perrin

\begingroup

\begin{frame}{Linéarisabilité}

  \begin{block}{Définition}
    Un objet partagé est linéarisable si : \\
    le résultat observable des opérations sur l'objet partagé d'une exécution est le même que celui d'un entrelacement de l'ordre ``happened before''
    \alert{tel que, si une opération $e$ se termine avant qu'une autre opération $e'$ ne commence, alors $e$ doit précéder $e'$ dans l'entrelacement}
  \end{block}

  \begin{exampleblock}{Exemple}
    
    \begin{center}
      \begin{tikzpicture}[y=7mm]
        \draw[process] (0,1) node[left]{$T_1$} to (8,1);
        \draw[process] (0,0) node[left]{$T_2$} to (8,0);
        \node[example, operation] at (3,0) {$increment()$ };
        \node[example, operation] at (6,1) {$get()$};
      \end{tikzpicture}
    \end{center}

    Un seul ordre de linéarisation possible : 
    \begin{itemize}
    \item $increment() ;~ get()$
      \begin{itemize}
      \item La lecture doit retourner 1
      \end{itemize}
    \end{itemize}
  \end{exampleblock}

\end{frame}

\endgroup
\endinput
